\subsection{Aproximações utilizadas na dinâmica}
\label{sec:aprox_dinamica}
A formulação apresentada inclui todos os termos da dinâmica relativa no referencial LVLH, o mesmo vale para as análises, o controle e o planejamento de trajetória. 
As principais hipóteses utilizadas neste trabalho são:
\begin{enumerate}[label=\roman*.]
    \item \textbf{Órbitas coplanares:} As trajetórias do \emph{chaser} e do \emph{target} possuem o mesmo plano orbital, eliminando termos de inclinação relativa e simplificando o vetor relativo.
    
    \item \textbf{Massas concentradas:} Ambas as espaçonaves são modelados como massas pontuais para fins de calculos orbitais, sendo desprezados efeitos de distribuição de massa e assimetrias gravitacionais na dinâmica translacional.
    
    \item \textbf{Pequena separação relativa:} Assume-se que o vetor $\vec s$ é pequeno em comparação ao raio orbital $r_t$, hipótese que permite aplicar expansões de Taylor, conforme \eqref{eq:linearizandoJacobiana}, em torno da posição do alvo e justifica as linearizações subsequentes.
\end{enumerate}

As aproximações realizadas permitem simplificações posteriores para a dedução das equações lineares de movimento relativo, apresentadas em seção posterior.